% ===================================================================
%  TABEL Nr. 4 — Volwassenheid en ouderdom.
%  Meme regime que les precedents.
%
%  LA PARENTE PAR ALLIANCE SE FAIT EN « skoon- » : skoonvader,
%  skoonseun, skoondogter, skoonsuster, swaer. L'ido la fabrique avec
%  « bo- » et le francais avec « beau- » ; les trois systemes se
%  repondent mot pour mot, ce qui est rare et qu'il faut saisir tant
%  qu'on l'a — chacun de ces mots porte un renvoi.
%
%  ET « ouma » EST LA GRAND-MERE, « oupa » LE GRAND-PERE. Les formes
%  neerlandaises « grootmoeder » et « grootvader » existent en
%  afrikaans mais s'y sentent ecrites ; l'usage dit ouma et oupa,
%  jusque dans les livres.
% ===================================================================

%%K t04-apar-1 apar
\VUsaut{4.0mm}
\VUtitre{15.0pt}{60}{TABEL Nr. 4}
\VUsaut{1.6mm}
\VUtitre{11.4pt}{0}{Volwassenheid en ouderdom.}
\VUsaut{3.2mm}

%%K t04-tit-1 sub
\VUcentre{11.4pt}{0}{\textit{Eerste tafereel.}}
\VUsaut{1.6mm}
\VUtitre{11.4pt}{0}{Die bruilofsmaal.}
\VUsaut{2.6mm}

%%K t04-tit-2 sub
\VUpk{10.2pt}{40}{Die herfs.}
\VUsaut{3.0mm}

%%K t04-c1-01-1 p
1. --- Die jongmense het groter geword. Een van hulle, Jan, trou. Ons is by die
\VUgras{bruilofsmaal} wat om dieselfde tafel die gesinne van die egpaar bymekaarbring.

%%K t04-c1-02-1 p
2. --- Ons is in die \VUgras{herfs}, in die maand \VUgras{September}. Die koue wind van
\VUgras{Oktober} en van \VUgras{November} het die platteland nog nie van sy groen
\VUgras{blaredak} beroof nie. Die aandlug is lou en geurig; daarom vind die ete onder die
\VUgras{veranda} \textsuperscript{(8)} by die tuin plaas.

%%K t04-c1-03-1 p
3. --- Ver weg, op die \VUgras{heuwel} \textsuperscript{(3)}, staan die huis van Jan se
ouers, 'n elegante \VUgras{kasteel} \textsuperscript{(1)} wat in die groen weggesteek is;
mooi wingerde, wat uitstekende wyn lewer, bedek die helling van die heuweltjie. Die
wynboer bewerk en versorg hulle.

%%K t04-c1-04-1 p
4. --- Bo die wingerde vlieg 'n vlug \VUgras{patryse} \textsuperscript{(2)} verby wat
deur die nadering van die \VUgras{veldwagter} \textsuperscript{(5)} verskrik is. Hierdie
een, met sy \VUgras{geweer} onder die arm en 'n \VUgras{wildsak} oor die skouer,
deursoek die \VUgras{bosse} \textsuperscript{(4)} met sy \VUgras{hond}
\textsuperscript{(6)}. Hy hou toesig oor die landgoed en keer dat die wildstropers die
wild uitroei.

%%K t04-c1-05-1 p
5. --- Buite die \VUgras{veranda} klim 'n \VUgras{leiwingerd} op waaraan digte
\VUgras{druiwetrosse} \textsuperscript{(7)} hang.

%%K t04-c1-06-1 p
6. --- Die mooi herfsblomme wat die \VUgras{tafel} \textsuperscript{(16)} versier, is
\VUgras{krisante} \textsuperscript{(9)}; die \VUgras{vader} \textsuperscript{(10)} van
die bruid het een daarvan in die knoopsgat van sy manelpak gesit.

%%K t04-c1-07-1 p
7. --- Die maaltyd loop ten einde. Die \VUgras{bediende} \textsuperscript{(11)} begin
die nagereg-skottels bring en sal binnekort die verskillende dele van die eetgerei
wegneem, \VUgras{lepels} \textsuperscript{(14)}, \VUgras{vurke} \textsuperscript{(12)},
\VUgras{messe} \textsuperscript{(13)}, \VUgras{skottels} \textsuperscript{(15)} wat 'n
mens nie meer nodig het nie.

%%K t04-tit-3 sub
\VUpk{10.2pt}{40}{Die familie.}
\VUsaut{3.0mm}

%%K t04-c1-08-1 p
8. --- Almal lyk vrolik, en die glimlag waarmee die \VUgras{moeder}
\textsuperscript{(17)} van die bruidegom na haar kinders kyk, bewys haar geluk.

%%K t04-c1-09-1 p
9. --- Hierdie huwelik verenig twee ryk gesinne van die streek. Jan, die
\VUgras{eggenoot} \textsuperscript{(18)}, is die \VUgras{seun} van 'n groot grondeienaar
en hy word die \VUgras{skoonseun} van 'n bekwame nyweraar.

%%K t04-c1-10-1 p
10. --- Die \VUgras{egpaar} is jonk en tog was hulle lank reeds \VUgras{verloof}. Die
\VUgras{jong vrou} \textsuperscript{(19)}, Marta, is bekoorlik in haar \VUgras{wit rok}
wat met lemoenbloeisels versier is.

%%K t04-c1-11-1 p
11. --- Haar \VUgras{skoonvader} \textsuperscript{(20)} is baie gelukkig om so 'n
vriendelike \VUgras{skoondogter} te hê, en die \VUgras{skoonmoeder}
\textsuperscript{(21)} van Jan glimlag as sy haar \VUgras{dogter} so mooi sien.

%%K t04-c1-12-1 p
12. --- Aan die punt van die tafel sit die ander \VUgras{genooides}
\textsuperscript{(22)}, \VUgras{familie, vriende, neefs, niggies}. 'n
\VUgras{tafelhoof} \textsuperscript{(23)}, met 'n servet onder die arm, bedien hulle
ywerig.

%%K t04-tit-4 sub
\VUpk{10.2pt}{40}{Die vriende.}
\VUsaut{3.0mm}

%%K t04-c1-13-1 p
13. --- Aan die regterkant van die tafel is 'n \VUgras{jongmeisie}
\textsuperscript{(24)}, 'n \VUgras{vriendin} van die jong bruid, dan die
\VUgras{broer} van hierdie een, wat haar \VUgras{strooijonker} \textsuperscript{(25)}
is. Hy gesels met sy \VUgras{strooimeisie} \textsuperscript{(26)}, die suster van die
bruidegom, wat deur hierdie huwelik die \VUgras{skoonsuster} van Marta word. Sy is 'n
bekoorlike jongmens. Sy het mooi \VUgras{juwele}, 'n \VUgras{pêrelhalssnoer} en 'n
goue \VUgras{armband}.

%%K t04-c1-14-1 p
14. --- Naby hulle is die meneer wat staan en sy glas oplig 'n \VUgras{vriend}
\textsuperscript{(27)} van die gesin. Hy is een van die \VUgras{getuies van die
bruidegom}. Hy stel 'n \VUgras{heildronk} in, drink op die \VUgras{gesondheid} van die
jong egpaar en wens hulle lange geluk toe. Hy is in 'n plegtige pak geklee, in 'n
manelpak, met \VUgras{gelakte skoene} \textsuperscript{(28)}. Hy is die metgesel van
die \VUgras{ouma} \textsuperscript{(29)} wat 'n \VUgras{weduwee} is.

%%K t04-c1-15-1 p
15. --- Die jongman in 'n \VUgras{manelpak} \textsuperscript{(31)}, met 'n dun swart
snor, is die \VUgras{swaer} \textsuperscript{(30)} van Jan. Hy is 'n voortreflike
ingenieur. Hy sit langs 'n baie elegante \VUgras{jong dame} \textsuperscript{(33)}, die
\VUgras{ouer suster} van Marta en die moeder van die fyn dogtertjie wat kom, terwyl sy
haar arm aan haar neef gee, om 'n blom aan te bied en haar 'n \VUgras{goeienaand} te
\VUgras{wens}. Daardie dame het baie mooi \VUgras{oorringe} en 'n elegante
\VUgras{hoofbedekking}.

%%K t04-c1-16-1 p
16. --- Die klein neef, so vreemd met sy \VUgras{bloes} \textsuperscript{(36)} en sy
bofferige \VUgras{broekie} \textsuperscript{(37)}, is baie beklaenswaardig. Hy is 'n
\VUgras{weeskind} \textsuperscript{(35)}. Baie jonk het hy sy vader en sy moeder verloor.
Sy oom is sy \VUgras{voog} \textsuperscript{(38)} en het ongetroud gebly om die arme kind
op te voed. Hy is 'n goedhartige man. Hy is 'n bietjie kaal en baie kortsigtig; hy
gebruik 'n \VUgras{knypbril}.

%%K t04-tit-5 sub
\VUsaut{2.4mm}
\VUpk{10.2pt}{40}{Die nagereg.}
\VUsaut{3.0mm}

%%K t04-c1-17-1 p
17. --- Agter die gaste, op 'n tafel wat met 'n \VUgras{tafeldoek} bedek is, is die
\VUgras{nagereg} \textsuperscript{(39)} voorberei. In 'n groot skaal is daar vrugte,
\VUgras{perskes} \textsuperscript{(40)}, \VUgras{pere} \textsuperscript{(41)},
\VUgras{appels} \textsuperscript{(42)}, \VUgras{appelkose} \textsuperscript{(43)} en
\VUgras{druiwe} \textsuperscript{(44)}. Naby, \VUgras{pynappels}
\textsuperscript{(45)}, 'n mooi \VUgras{koek} \textsuperscript{(46)},
\VUgras{likeurbottels} \textsuperscript{(48)} en \VUgras{sjampanje}
\textsuperscript{(49)} en ook stapels skoon \VUgras{borde} \textsuperscript{(47)}.

%%K t04-c1-18-1 p
18. --- Die \VUgras{kelderman} \textsuperscript{(50)} trek 'n \VUgras{bottel}
\textsuperscript{(52)} ou wyn oop met 'n \VUgras{kurktrekker} \textsuperscript{(51)}.
Die \VUgras{kurkprop} \textsuperscript{(53)} lyk baie moeilik om uit te trek. Daardie
kelderman het 'n \VUgras{servet} \textsuperscript{(54)} onder sy arm.

%%K t04-c1-19-1 p
19. --- Na die ete sal die here na die rookkamer gaan, en die dames sal hulle vir die
bal gaan gereedmaak.

%%K t04-tit-6 sub
\VUsaut{5.0mm}
\VUcentre{11.4pt}{0}{\textit{Tweede tafereel.}}
\VUsaut{1.6mm}
\VUpk{11.4pt}{40}{Die oupa se verjaardag.}
\VUsaut{2.6mm}

%%K t04-tit-7 sub
\VUpk{10.2pt}{40}{Die besoek.}
\VUsaut{3.0mm}

%%K t04-c2-01-1 p
1. --- Tien jaar het verbygegaan, Jan se ouers het oud geword en is baie lief vir hulle
drie kleinkinders, twee \VUgras{kleinseuns} \textsuperscript{(2)} en een
\VUgras{kleindogter} \textsuperscript{(3)}. Omdat dit vandag \VUgras{die oupa se
verjaardag} is, kom die kinders hulle 'n gelukkige naamdag toewens.

%%K t04-c2-02-1 p
2. --- Die klein Pieter is nog baie jonk, hy is net drie jaar oud. Hy sien die
\VUgras{kat} \textsuperscript{(1)} nader kom. Hy wil sy mooi sagte \VUgras{pels} en sy
lang \VUgras{stert} streel, hy dink nie daaraan dat onder die sierlike \VUgras{pote}
kwaadaardige skerp naels wegkruip wat hom miskien sal krap nie.

%%K t04-c2-03-1 p
3. --- Sy sussie Gabriela, wat hom die hand gee, is veel ernstiger, en Marcellus met sy
groot \VUgras{jas} \textsuperscript{(4)} en leer \VUgras{gordel} \textsuperscript{(5)}
lyk 'n bietjie manlik geword, ondanks sy kort broekie wat sy \VUgras{kouse}
\textsuperscript{(6)} laat sien.

%%K t04-c2-04-1 p
4. --- Hulle vader het sy dik winter\VUgras{oorjas} \textsuperscript{(7)} met 'n
\VUgras{pelskraag} \textsuperscript{(8)} aangetrek en in sy \VUgras{sak}
\textsuperscript{(9)} het hy 'n serp. Sy vrou het haar vilthoed wat met \VUgras{vere}
\textsuperscript{(10)} versier is, haar \VUgras{jakkie} \textsuperscript{(11)}, haar
\VUgras{boa} \textsuperscript{(12)} en haar \VUgras{mof} \textsuperscript{(13)}.

%%K t04-c2-05-1 p
5. --- Dit is baie koud, want die winter heers. Die \VUgras{sneeu} \textsuperscript{(19)}
het die hele dag geval en die dakke van die huise is heeltemal wit. Saam met die
\VUgras{nag} word die koue nog skerper. In die hemel skyn die \VUgras{maan}
\textsuperscript{(17)} en die \VUgras{sterre} \textsuperscript{(18)} en dring deur
\VUgras{die duisternis}.

%%K t04-tit-8 sub
\VUsaut{4.2mm}
\VUpk{10.2pt}{40}{Die siekte.}
\VUsaut{3.0mm}

%%K t04-c2-06-1 p
6. --- Die arme \VUgras{blinde} \textsuperscript{(20)} wat die plein voor die
\VUgras{apteek} \textsuperscript{(16)} oorsteek, gelei deur sy \VUgras{poedel}
\textsuperscript{(21)}, is baie ongelukkig. Justina, die \VUgras{diensmeisie}
\textsuperscript{(14)}, bring uit die kombuis 'n \VUgras{bak} \textsuperscript{(15)}
baie warm \VUgras{kruietee} wat mildelik gesoet is.

%%K t04-c2-07-1 p
7. --- Die \VUgras{ouma} \textsuperscript{(23)}, wat haar \VUgras{bril}
\textsuperscript{(26)} op die tafel wil gaan soek om die \VUgras{voorskrif}
\textsuperscript{(27)} te lees wat die \VUgras{dokter} \textsuperscript{(25)} pas
geskryf het, staan uit haar leunstoel op en steun op haar \VUgras{kierie}
\textsuperscript{(24)}.

%%K t04-c2-08-1 p
8. --- Dikwels ly sy aan \VUgras{ongesteldhede, duiselings, verkoue, tandpyne}, haar
bene is swak en haar oë is nie meer baie goed nie. Ten spyte daarvan spot sy graag met
haar ou \VUgras{dokter} wat haar altyd 'n \VUgras{drankie} \textsuperscript{(28)} wil
voorskryf. Vandag is sy baie bly, want sy het haar jong gesin om haar.

%%K t04-c2-09-1 p
9. --- Dit is gesellig in die goed toegemaakte kamer. In die marmer \VUgras{kaggel}
\textsuperscript{(29)} vlam 'n \VUgras{pragtige vuur} \textsuperscript{(30)}, en 'n mens
voel gelukkig in die sagte \VUgras{lig} van die \VUgras{lamp} \textsuperscript{(31)} wat
sopas aangesteek is.

%%K t04-tit-9 sub
\VUsaut{4.2mm}
\VUpk{10.2pt}{40}{Die oupa.}
\VUsaut{3.0mm}

%%K t04-c2-10-1 p
10. --- Selfs die \VUgras{oupa} \textsuperscript{(33)} het vergeet dat hy siek was, dat
sy \VUgras{rumatiek} hom aan sy \VUgras{leunstoel} \textsuperscript{(34)} vasmaak en dat
hy gister nog \VUgras{koors en floutes} gehad het. Hy onthou nie meer dat hy die laaste
winter aan \VUgras{longontsteking} gely het en dat 'n hees en pynlike \VUgras{hoes} hom
baie laat ly het nie.

%%K t04-c2-10-2 p
En tog is hy baie moeilik gesond gemaak. Van nou af gaan dit 'n bietjie beter met hom.
Maar sy been, wat hy op 'n \VUgras{kussing} \textsuperscript{(38)} laat rus wat op 'n
klein \VUgras{bankie} \textsuperscript{(39)} gesit is, maak hom ook seer.

%%K t04-c2-11-1 p
11. --- Wanneer hy sorgvuldig in sy warm \VUgras{kamerjas} \textsuperscript{(35)} met
dik perlemoen\VUgras{knope} \textsuperscript{(36)} toegedraai is en wanneer hy naby hom,
om die koerant vir hom te lees, die klein \VUgras{weesdogtertjie} \textsuperscript{(40)}
het wat in 'n wit \VUgras{kappie}, in 'n blou \VUgras{rokkie} \textsuperscript{(42)} en
'n \VUgras{skouermanteltjie} \textsuperscript{(41)} geklee is, kla hy nie.

%%K t04-c2-12-1 p
12. --- Elke jaar, wanneer die \VUgras{kalender} \textsuperscript{(43)} weer die
\VUgras{datum} van die eerste \VUgras{Desember} aandui, wag hy ongeduldig op sy
\VUgras{kleinkinders}, want hy weet dat hulle daardie belangrike \VUgras{datum} nie sal
vergeet nie.

%%K t04-c2-13-1 p
13. --- Met ontroering onthou hy die tyd baie lank gelede toe hy self die roemryke
keiserlike \VUgras{maarskalk} 'n gelukkige naamdag gaan toewens het, wie se
\VUgras{portret} \textsuperscript{(44)} teen die muur hang en wat die
\VUgras{oupagrootjie} van sy kinders is.
\VUsaut{6.0mm}
\VUfilet{22mm}
